\documentclass[12pt]{article}
\usepackage[utf8]{inputenc}
\usepackage[T1]{fontenc}
\usepackage[margin=1in]{geometry}
\usepackage{setspace}
\usepackage{xcolor}
\usepackage[colorlinks=true, urlcolor=blue!60!black, linkcolor=blue!60!black]{hyperref}
\setlength{\parindent}{2em}
\setlength{\parskip}{1.25em}
\renewcommand{\baselinestretch}{1.0}
\begin{document}
\begin{center}
{\LARGE \textit{Rucho v. Common Cause}, 588 U.S. 684 (2019)\par}
\vspace{0.8em}
{\large Case Brief --- Four-Page Edition\par}
\end{center}
\vspace{1em}

\section{Memory Jogger}

Partisan gerrymandering claims are nonjusticiable political questions beyond the federal courts' jurisdiction, because the Constitution supplies no clear, manageable, and politically neutral standard for deciding how much partisanship in districting is too much (\href{https://www.courtlistener.com/opinion/4633469/rucho-v-common-cause/}{CourtListener}).

\section{Detailed Case Facts}

\textbf{North Carolina.} After two districts were struck down as racial gerrymanders, the General Assembly, controlled by Republicans, adopted a 2016 congressional plan with ``Partisan Advantage'' as an express criterion. Representative David Lewis explained that he drew 10 Republican and 3 Democratic seats only because he did not believe an 11--2 map was possible. Republicans won 10 seats in 2016 with 53\% of the statewide vote and 9 of 12 decided seats in 2018 with 50\%.

\textbf{Maryland.} Democrats, led by Governor Martin O'Malley and Congressman Steny Hoyer, who called himself a ``serial gerrymanderer,'' redrew the map in 2011 to flip the Sixth District. Although equal population required moving about 10,000 residents, the plan moved roughly 360,000 out and 350,000 in, cutting registered Republicans by about 66,000. Democrats won 7 of 8 seats in every election from 2012 to 2018 without exceeding 65\% of the vote.

\section{Procedural History}

In North Carolina, a three-judge court struck down the 2016 plan under the Equal Protection Clause, Article I, and (2--1) the First Amendment. The Supreme Court vacated in light of \textit{Gill v. Whitford}; on remand the panel found standing and invalidated 12 of 13 districts. In Maryland, a three-judge court granted summary judgment to Republican voters on First Amendment grounds and enjoined the map. On direct appeal, the Supreme Court consolidated the cases and, on June 27, 2019, vacated both judgments with instructions to dismiss for lack of jurisdiction.

\section{Judicial Votes}

Vacated and remanded, \textbf{5--4}.
\begin{itemize}
\item \textbf{Majority:} Chief Justice Roberts, joined by Justices Thomas, Alito, Gorsuch, and Kavanaugh.
\item \textbf{Concurrences:} None.
\item \textbf{Dissent:} Justice Kagan, joined by Justices Ginsburg, Breyer, and Sotomayor.
\end{itemize}

\section{Holding}

Partisan gerrymandering claims---under the Equal Protection Clause, the First Amendment, Article I, \S~2, or the Elections Clause---present nonjusticiable political questions because there are no judicially discoverable and manageable standards for deciding when partisan consideration becomes unconstitutional. The holding is federal and jurisdictional. The Court's statements that excessive partisanship is ``incompatible with democratic principles,'' and its survey of state constitutions, commissions, state courts, and congressional bills, are dicta.

\section{Analysis of Opinions}

\textbf{Majority (Roberts).} Partisan gerrymandering predates the Constitution, from Patrick Henry's effort against Madison to Elbridge Gerry's 1812 map, and the Framers assigned districting to state legislatures checked by Congress. One-person, one-vote claims are administrable by arithmetic and racial gerrymanders target a suspect classification, but ``a jurisdiction may engage in constitutional political gerrymandering.'' The question is therefore one of degree, and the claims rest on an unstated demand for proportional representation, which the Constitution does not require. Defining fairness is political, not legal. The Court rejected the North Carolina intent-and-effects test (judges cannot ``beat the point spread''), the First Amendment theory (maps restrict no speech), the dissent's ``extreme outlier'' baseline (arbitrary and still unanswered as to ``too much''), and Article I claims. Reform lies with States and Congress.

\textbf{Kagan (dissenting).} Extreme gerrymanders let politicians choose their voters, violating equal protection through packing and cracking and the First Amendment by penalizing political views, and big data makes them more durable than ever. Lower courts had converged on a neutral test---predominant intent to entrench, substantial dilutive effect, and no legitimate justification---using thousands of computer-simulated maps built from each State's own criteria to identify ``extreme outliers.'' To ``how much is too much?'' she answered, ``This much is too much.'' Politicians will not dismantle their own gerrymanders, and fewer than half the States allow initiatives.

\section{Examples: Future Applications}

\begin{enumerate}
\item \textbf{Same-Side (Nonjusticiable): State Legislative Map.} Software-drawn state senate districts that lock in a 65\% supermajority on 48\% of the vote must be dismissed in federal court.
\item \textbf{Same-Side (Nonjusticiable): Mid-Decade Redistricting.} A new majority's redraw to oust two incumbents, with exact population equality and no racial classifications, is barred.
\item \textbf{Opposite-Side (Justiciable): State Constitutions.} State courts may enforce ``Free and Equal Elections'' clauses, and \textit{Moore v. Harper} (2023) confirmed that the Elections Clause does not insulate legislatures from state judicial review.
\item \textbf{Opposite-Side (Justiciable): Race and Population.} A 12\% population deviation or race-predominant packing of Black voters remains reviewable.
\item \textbf{Fence-Sitter: Race and Party Intertwined.} Where race and party correlate, moving Black Democratic voters for claimed partisan reasons forces plaintiffs to disentangle the two, as \textit{Alexander v. South Carolina NAACP} (2024) demanded.
\end{enumerate}

\section{Critique}

\textbf{Originalist views.} Conservatives praised the deference to the Elections Clause's assignment of oversight to Congress. John McGinnis called it ``silent originalism'': the Equal Protection Clause, as an original matter, created no right against partisan dilution, so the political question label was unnecessary. Other formalists question whether the label is coherent when the real issue is whether any right exists.

\textbf{Progressive views.} Election-law scholars call \textit{Rucho} an abdication where political self-correction fails because incumbents control the rules. The \textit{Harvard Law Review} noted that its prudential reasoning leaves state courts free to hear claims federal courts cannot, and critics say it greenlights technological entrenchment and invites recasting partisan claims as racial ones.

\section{Key Quotations}

\begin{quote}
``Excessive partisanship in districting leads to results that reasonably seem unjust. But the fact that such gerrymandering is `incompatible with democratic principles,' \textit{Arizona State Legislature}, 576 U.\,S., at 791 (brackets omitted), does not mean that the solution lies with the Federal Judiciary. We conclude that partisan gerrymandering claims present political questions beyond the reach of the federal courts.''
\end{quote}
\noindent --- Chief Justice Roberts, for the Court, 588 U.S. at 718.

\begin{quote}
``No one can accuse this Court of having a crabbed view of the reach of its competence. But we have no commission to allocate political power and influence in the absence of a constitutional directive or legal standards to guide us in the exercise of such authority. `It is emphatically the province and duty of the judicial department to say what the law is.' \textit{Marbury v. Madison}, 1 Cranch, at 177. In this rare circumstance, that means our duty is to say `this is not law.' ''
\end{quote}
\noindent --- Chief Justice Roberts, for the Court, 588 U.S. at 721.

\begin{quote}
``For the first time ever, this Court refuses to remedy a constitutional violation because it thinks the task beyond judicial capabilities.''
\end{quote}
\noindent --- Justice Kagan, dissenting, 588 U.S. at 721.

\end{document}
